% GENERATED FILE. Edit canonical lessons, then run npm run generate:books.
% canonical-source-hash: fnv1a64:ce0c67200976814c
% canonical-lessons: GE-C412-bahnsteig, GE-C412-durchsage, GE-C412-abteil, GE-C412-schaffner, GE-C412-kontrollieren

\chapter{Auf dem Bahnsteig}
\label{ch:ge-auf-dem-bahnsteig}
\hlchaptermodality{412}

\begin{chapteropening}
\textbf{By the end of this chapter:} \emph{I can manage on a platform: the platform, the announcement, the compartment, the conductor and checking tickets.}

Complete the last lesson of chapter 412: I can manage on a platform: the platform, the announcement, the compartment, the conductor and checking tickets.
\end{chapteropening}

\section[der Bahnsteig]{der Bahnsteig — the platform (at a station)}

\label{lesson:GE-C412-bahnsteig}



\begin{quote}
Before the new one: say the German for to stamp, then the German for to get a ticket.
\end{quote}

\subsection*{You'll want to know: der Bahnsteig}
\textbf{der Bahnsteig} — \textquotedblleft{}the platform (at a station)\textquotedblright{}.

\begin{quote}
Unser Zug fährt heute von einem anderen Bahnsteig.
\end{quote}

\begin{quote}
\emph{Our train leaves from a different platform today.}
\end{quote}

Plural: \emph{die Bahnsteige}. The track is \emph{das Gleis}. In Switzerland: \emph{der Perron}.

\begin{cousinweb}
\emph{Bahnsteig} is built from \emph{Bahn} 'track, railway' and \emph{Steig} 'steep path', from \emph{steigen} 'to climb', inherited from Old High German \emph{stīgan} (EtymDB).
\end{cousinweb}

\begin{grammarlens}[title={What you've built}]
One more word for telling, explaining and describing at B1.
\end{grammarlens}

\subsection*{Your turn}
\begin{itemize}
\raggedright
  \item \emph{Say it:} \emph{der Bahnsteig}
  \item \emph{Say it:} \emph{der Bahnsteig}, once more
  \item \emph{Say it:} a sentence of your own with \emph{der Bahnsteig}
  \item \emph{Recall:} say the German for to stamp, then the German for to get a ticket, then say \emph{der Bahnsteig} again
\end{itemize}

\begin{tcolorbox}[breakable,colback=teal!4,colframe=teal!35!black,title={Before you move on}]
What does der Bahnsteig mean? (\textquotedblleft{}The platform (at a station)\textquotedblright{}.) Say it once more.
\end{tcolorbox}
\section[die Durchsage]{die Durchsage — the announcement (over a loudspeaker)}

\label{lesson:GE-C412-durchsage}



\begin{quote}
Before the new one: say the German for to get a ticket, then the German for the platform.
\end{quote}

\subsection*{You'll want to know: die Durchsage}
\textbf{die Durchsage} — \textquotedblleft{}the announcement (over a loudspeaker)\textquotedblright{}.

\begin{quote}
Laut der Durchsage hat der Zug zehn Minuten Verspätung.
\end{quote}

\begin{quote}
\emph{According to the announcement the train is ten minutes late.}
\end{quote}

Plural: \emph{die Durchsagen}.

\begin{cousinweb}
\emph{Durchsage} is built from \emph{durchsagen} 'to pass a message on', from \emph{durch} 'through' and \emph{sagen} 'to say', inherited from Proto-Germanic \emph{sagjaną} (EtymDB).
\end{cousinweb}

\begin{grammarlens}[title={What you've built}]
One more word for telling, explaining and describing at B1.
\end{grammarlens}

\subsection*{Your turn}
\begin{itemize}
\raggedright
  \item \emph{Say it:} \emph{die Durchsage}
  \item \emph{Say it:} \emph{die Durchsage}, once more
  \item \emph{Say it:} a sentence of your own with \emph{die Durchsage}
  \item \emph{Recall:} say the German for to get a ticket, then the German for the platform, then say \emph{die Durchsage} again
\end{itemize}

\begin{tcolorbox}[breakable,colback=teal!4,colframe=teal!35!black,title={Before you move on}]
What does die Durchsage mean? (\textquotedblleft{}The announcement (over a loudspeaker)\textquotedblright{}.) Say it once more.
\end{tcolorbox}
\section[das Abteil]{das Abteil — the compartment (on a train)}

\label{lesson:GE-C412-abteil}



\begin{quote}
Before the new one: say the German for the platform, then the German for the announcement.
\end{quote}

\subsection*{You'll want to know: das Abteil}
\textbf{das Abteil} — \textquotedblleft{}the compartment (on a train)\textquotedblright{}.

\begin{quote}
In unserem Abteil sitzen nur zwei Leute.
\end{quote}

\begin{quote}
\emph{Only two people are sitting in our compartment.}
\end{quote}

Plural: \emph{die Abteile}. Many new trains have open carriages instead.

\begin{cousinweb}
\emph{Abteil} is built on \emph{abteilen} 'to divide off', from \emph{Teil} 'part', inherited from Old High German \emph{teil} and Proto-Germanic \emph{dailą} (EtymDB).
\end{cousinweb}

\begin{grammarlens}[title={What you've built}]
One more word for telling, explaining and describing at B1.
\end{grammarlens}

\subsection*{Your turn}
\begin{itemize}
\raggedright
  \item \emph{Say it:} \emph{das Abteil}
  \item \emph{Say it:} \emph{das Abteil}, once more
  \item \emph{Say it:} a sentence of your own with \emph{das Abteil}
  \item \emph{Recall:} say the German for the platform, then the German for the announcement, then say \emph{das Abteil} again
\end{itemize}

\begin{tcolorbox}[breakable,colback=teal!4,colframe=teal!35!black,title={Before you move on}]
What does das Abteil mean? (\textquotedblleft{}The compartment (on a train)\textquotedblright{}.) Say it once more.
\end{tcolorbox}
\section[der Schaffner]{der Schaffner — the conductor, the ticket inspector (on a train)}

\label{lesson:GE-C412-schaffner}



\begin{quote}
Before the new one: say the German for the announcement, then the German for the compartment.
\end{quote}

\subsection*{You'll want to know: der Schaffner}
\textbf{der Schaffner} — \textquotedblleft{}the conductor, the ticket inspector (on a train)\textquotedblright{}.

\begin{quote}
Der Schaffner kommt und will unsere Fahrkarten sehen.
\end{quote}

\begin{quote}
\emph{The conductor comes and wants to see our tickets.}
\end{quote}

Same form in the plural. A woman: \emph{die Schaffnerin}. Today the official word is \emph{der Zugbegleiter}.

\begin{cousinweb}
\emph{Schaffner} is built on \emph{schaffen} in the old sense 'to arrange, to manage', inherited from Middle High German \emph{schaffen} (EtymDB); the word is usually said to have first named a steward who ran a household.
\end{cousinweb}

\begin{grammarlens}[title={What you've built}]
One more word for telling, explaining and describing at B1.
\end{grammarlens}

\subsection*{Your turn}
\begin{itemize}
\raggedright
  \item \emph{Say it:} \emph{der Schaffner}
  \item \emph{Say it:} \emph{der Schaffner}, once more
  \item \emph{Say it:} a sentence of your own with \emph{der Schaffner}
  \item \emph{Recall:} say the German for the announcement, then the German for the compartment, then say \emph{der Schaffner} again
\end{itemize}

\begin{tcolorbox}[breakable,colback=teal!4,colframe=teal!35!black,title={Before you move on}]
What does der Schaffner mean? (\textquotedblleft{}The conductor, the ticket inspector (on a train)\textquotedblright{}.) Say it once more.
\end{tcolorbox}
\section[kontrollieren]{kontrollieren — to check, to inspect; to control}

\label{lesson:GE-C412-kontrollieren}



\begin{quote}
Before the new one: say the German for the compartment, then the German for the conductor.
\end{quote}

\subsection*{You'll want to know: kontrollieren}
\textbf{kontrollieren} — \textquotedblleft{}to check, to inspect; to control\textquotedblright{}.

\begin{quote}
Im Zug kontrolliert der Schaffner die Fahrkarten.
\end{quote}

\begin{quote}
\emph{On the train the conductor checks the tickets.}
\end{quote}

Past participle without \emph{ge-}: \emph{er hat die Pässe kontrolliert}. \emph{Die Kontrolle} — the check.

\begin{cousinweb}
\emph{kontrollieren} was borrowed from French \emph{contrôler} (EtymDB), built on Old French \emph{contrerole} 'counter-roll', a second list kept to check the first.
\end{cousinweb}

\begin{grammarlens}[title={What you've built}]
One more word for telling, explaining and describing at B1.
\end{grammarlens}

\subsection*{Your turn}
\begin{itemize}
\raggedright
  \item \emph{Say it:} \emph{kontrollieren}
  \item \emph{Say it:} \emph{kontrollieren}, once more
  \item \emph{Say it:} a sentence of your own with \emph{kontrollieren}
  \item \emph{Recall:} say the German for the compartment, then the German for the conductor, then say \emph{kontrollieren} again
\end{itemize}

\begin{tcolorbox}[breakable,colback=teal!4,colframe=teal!35!black,title={Before you move on}]
What does kontrollieren mean? (\textquotedblleft{}To check, to inspect; to control\textquotedblright{}.) Say it once more.
\end{tcolorbox}
